\clearpage\label{answers}
{\sans\bfseries\LARGE\color{ink} Answers}\par
\small
Try every puzzle first. Where an answer shows a calculation, check the steps, not just the result.

\subhead{Chapter 1}
\textbf{Multix on the clip test.} Throwing one plane hard and the other softly is a second difference, so the test cannot say whether the clip or the push made the gap. \textbf{Spot the changes (Dev).} He meant to change plant food. Light also differs: sun against a dark cupboard. The test is not fair. \textbf{Keep box.} Kept: the paper, the spot, the thrower, the fold. Changed: only the one thing being tested.

\textbf{1.1} B is fair. A changes the surface and the kind of ball as well as the weight. C changes the room as well as the music. A has the most hidden differences.\\
\textbf{1.2} (a) C (b) M (c) K (d) K.\\
\textbf{1.3} Eli. If the amounts of water differ, that is a second change, and then the result cannot say which change mattered.\\
\textbf{1.4} (a) F. (b) U: time of day and tiredness differ. (c) F. (d) U: the torch and the toy differ. (e) F.\\
\textbf{1.5} (a) The brand of towel. (b) The amount of towel and the size of the container. (c) Use the same amount of each towel in the same container.\\
\textbf{1.6} (a) paper size. (b) water temperature. (c) the height of the tower.

\subhead{Chapter 2}
\textbf{Multix on his own test.} He threw one plane hard and the other softly, so the push differs as well as the clip. The 4.5~m gap cannot be blamed on the clip. \textbf{Keep box.} Kept: the typical distance and the conditions that stayed the same. Set aside: which throw was which.

\textbf{2.1} (a) $3.9+4.2+4.5=12.6$, average 4.2~m, spread $4.5-3.9=0.6$~m. (b) $5.1+5.4+5.7=16.2$, average 5.4~m, spread 0.6~m. (c) $5.4-4.2=1.2$~m.\\
\textbf{2.2} (a) Change the ramp length; measure how far the car rolls; keep the same car, floor and ramp tilt. (b) Change the water temperature; measure the minutes the ice takes to melt; keep the same ice cube, amount of water and starting time.\\
\textbf{2.3} Any fair items, such as the same paper size, the same fold, the same thrower, the same launch line.\\
\textbf{2.4} Sam. Three throws each is steadier than one. One throw can be a fluke.

\subhead{Chapter 3}
\textbf{Multix on his card.} No. ``A good way'' is vague, so two people will do different tests. \textbf{Keep box.} Kept: the question, the one change, what to measure, what to keep the same and how many times. Left off: opinions and stories about the throws.

\textbf{3.1} Examples: how the car is released (let go or pushed), the ramp's steepness, or the ramp's surface. A hidden difference there changes the result.\\
\textbf{3.2} Throw the plain plane, then the clip plane: 4. Fold two planes: 2. Find the averages: 6. Write down the distance: 5. Put a clip on one plane: 3. Tape a launch line: 1.\\
\textbf{3.3} Examples: (a) ``Use a cup that holds 250~mL.'' (b) ``Wait 10 minutes, timed with a clock.'' (c) ``Stir 20 times, one turn a second.''\\
\textbf{3.4} Examples: a different push, a different launch height, a different kind of paper, or measuring to the first landing in one class and to where the plane stopped in the other.\\
\textbf{3.5} Question: does a higher ramp send the car farther? Change: the number of books under the ramp. Measure: how far the car rolls past the ramp. Keep the same: the car, the ramp, the floor and letting go without a push. Times: three rolls each, in a mixed order.\\
\textbf{3.6} (a) The person who got 6.9. (b) Examples: How hard did you throw? From where? Was the same paper used? How did you measure?\\
\textbf{Second card.} For example: ``Use one layer of each wrapping, about the same thickness, or record the thickness.''

\subhead{Chapter 4}
\textbf{Multix on flaps.} A gap of 0.1~m is far smaller than the spreads of 0.9 and 1.1~m, so the test cannot tell. \textbf{Keep box.} Kept: the typical distance and the spread. Set aside: which throw was which.

\textbf{4.1} X: total 16.0, average 3.2, spread 0.4. Y: total 20.0, average 4.0, spread 0.4. Gap $4.0-3.2=0.8$~m. Yes, the gap is bigger than both spreads.\\
\textbf{4.2} No. The averages match, but each spread is about 1.5~m, so the groups overlap heavily. Throw more before you decide.\\
\textbf{4.3} (a) Plain $12.8\div3\approx4.3$, clip $13.1\div3\approx4.4$. (b) Plain spread 0.5, clip spread 0.7. (c) Throw more: the gap, 0.1~m, is much smaller than the spreads.\\
\textbf{4.4} Examples: the thrower may have warmed up, got tired or got better, or the wind may have changed. A mixed order spreads such changes across both planes.\\
\textbf{4.5} (a) A: $210\div5=42$ cm, B: $265\div5=53$ cm. (b) Both spreads are 4. (c) Gap 11~cm, which is bigger than both spreads.

\subhead{Chapter 5}
\textbf{Multix on wool and metal.} No. The sock is also much thicker than the foil, so material and thickness are both different. \textbf{Keep box.} Kept: the ice, the saucers, the spot and the start time. We also let the thickness change by accident.

\textbf{5.1} (a) Foil $(44+48)\div2=46$~min, paper towel $(60+62)\div2=61$~min. (b) Gap 15~min. (c) Foil spread 4, paper towel spread 2. (d) Yes, 15 is bigger than 4 and 2.\\
\textbf{5.2} (a) The sponges are different sizes. (b) Use sponges of the same size, soaked for the same time. (c) $62-10=52$~g.\\
\textbf{5.3} B and C. In B the ice sizes differ. In C the sun and shade differ as well as the wrapping.\\
\textbf{5.4} Change: salt or no salt, with the same amount of salt each time. Measure: minutes until the ice has melted. Keep the same: ice size, saucer, spot and starting time.\\
\textbf{5.5} B. It names a real limit of the test. A is silly and C is vague.

\subhead{Chapter 6}
\textbf{Multix on one throw.} No. One throw of one plane cannot show anything about all clips and all planes. \textbf{Keep box.} See the recap table.

\textbf{6.1} (a) S (b) B (c) G (d) S.\\
\textbf{6.2} For example: ``In this test, planes with flaps flew about 4.3~m and flat planes about 4.2~m. The spreads were bigger than the gap, so the test cannot tell whether flaps matter.''\\
\textbf{6.3} Not fair. Sun against shade is a second change, as well as the plant food.\\\\
\textbf{6.4} (a) One throw each, so a fluke can decide. (b) The weights differ, so the tape is not the only change. (c) Tiredness grew during the second group, so the order was not mixed.\\
\textbf{6.5} (a) $3.8+4.0+4.2+4.0+4.0=20.0$, average 4.0. (b) $5.1+5.3+5.0+5.4+5.2=26.0$, average 5.2. (c) Gap 1.2. (d) Spreads 0.4 and 0.4.\\
\textbf{6.6} Class B. It stays inside the test and gives the number.
\textbf{Slop check.} (a) You cannot tell: there is no title, no labels and no units. (b) Many fonts and colours, stars, a pie shape and rotated filler text. (c) Names for the bars and the unit, minutes. (d) The pie, the stars and the jokes.\\
\textbf{Redraw.} One title, labelled bars and axis with minutes, one or two calm colours and nothing else.

\textbf{Whole-book table.} Chapter 2: simplified a set of throws to an average and a spread, kept the typical distance and the conditions, set aside which throw was which. Chapter 3: simplified a test to a card, kept everything a stranger needs, set aside the thrower's opinions. Chapter 4: simplified five throws to an average and spread, kept the gap compared with the spread, set aside the order of the throws. Chapter 5: simplified a melting cube to a time, kept the minutes and the wrapping, set aside the exact shape. Chapter 6: simplified a test to a report, kept the claim and its limits, set aside the detail of each step.

\clearpage\label{notes}
{\sans\bfseries\LARGE\color{ink} Words}\par
\small
\begin{multicols}{2}
\word{average}{the total of the readings divided by how many there are.}
\word{claim}{a sentence that says what a test showed, within the limits of the test.}
\word{evidence}{the numbers and observations a test gives you.}
\word{fair test}{a test where only one thing is changed and everything else is kept the same.}
\word{gap}{the difference between two averages.}
\word{keep the same}{to hold a condition fixed in every trial.}
\word{launch line}{a marked spot that every throw starts from.}
\word{repeat}{to run the same test again, to see whether the result comes back.}
\word{spread}{the biggest reading minus the smallest.}
\word{test card}{a written list of the question, the change, the measurement and everything kept the same.}
\word{unfair test}{a test where more than one thing differs between the versions being compared.}
\word{variable}{something in a test that can change.}
\end{multicols}

\vspace{0.4em}
{\sans\bfseries\large\color{sky!70!black} Notes for grown-ups}\par
This book teaches that a result is only as good as the test behind it. Each chapter has one activity at a table or on the floor, one worked example with made-up practice numbers (labelled as such), one trust box, one ``what did we keep?'' prompt, and a page of puzzles.

\textbf{Materials and safety.} Scrap paper, paper clips, a tape measure, a timer, ice cubes, saucers, aluminium foil, a paper towel and a sock. Throw planes gently, never toward faces or pets. Keep clips away from small children. Let an adult handle the freezer, and do not taste test materials.

\textbf{Check boxes.} The trust boxes say what a check shows and what it does not show. Repeating a test shows that it can be repeated, not that the explanation is right. A gap bigger than the spread is only a rule of thumb, not a proof.

\textbf{The practice numbers} in the worked examples (planes, ice, sponges) were made up so that the arithmetic is easy to check. They are not results of real experiments, and real results will differ. The ice example also hides a second change, the thickness of the wrapping, on purpose.

\subhead{Sources}
\begin{itemize}[leftmargin=1.3em,itemsep=0.2em]
\item R. A. Fisher. \emph{The Design of Experiments}. Oliver and Boyd, 1935. For changing one factor at a time and the logic of fair comparison.
\item J. Lind. \emph{A Treatise of the Scurvy}. Edinburgh, 1753. For the 1747 comparison of treatments.
\item J. R. Taylor. \emph{An Introduction to Error Analysis}, 2nd edition. University Science Books, 1997. For repeated measurement and spread.
\end{itemize}

\clearpage
\thispagestyle{blank}
\begin{tikzpicture}[remember picture,overlay]
  \fill[sun!35] ($(current page.south west)+(0.4in,0.4in)$) rectangle ($(current page.north east)-(0.4in,0.4in)$);
  \node[anchor=north,align=center,font=\sffamily\bfseries,text=ink,text width=5.4in] at ($(current page.north)+(0,-1.4in)$) {\Large Next in Science\\[6pt]\fontsize{28}{30}\selectfont Shadows That Tell Time};
  \node[align=center,font=\sffamily,text width=4.6in,text=ink] at (current page.center) {A stick in the sun draws a line on the ground, and the line moves. Book 3 shows how light travels, why shadows grow and shrink, and how to build a clock out of a stick, a stone and a sunny day.};
  \node at ($(current page.south)+(0,2.6in)$) {\scalebox{3}{\multixbot[ok]}};
  \node[font=\sffamily\bfseries,text=ink] at ($(current page.south)+(0,0.9in)$) {Ages 9 to 11 \quad$\star$\quad Science Book 2};
\end{tikzpicture}
